% ===================================================================
%  TABLEAU N° 14 — La rue. --- Les commerçants.
%
%  Traduction, faite en 2026, en francais standard du Quebec, sur
%  l'ido de texto/io. Regles de la colonne : voir 10-tableau-01.tex.
%
%  LES NEUF CHOIX PROPRES A CE TABLEAU :
%    (12) balayeuse mécanique ... voiture de balayage
%    (78) borne-fontaine ........ fontaine de pierre
%    (7)  liquoriste ............ marchand de spiritueux
%    (55) voyageur de commerce .. représentant de commerce
%    (92) marchande de fleurs ... fleuriste
%    (98) bonne d'enfant ........ gardienne d'enfants
%    (64) magasin de nouveautés . magasin de tissus
%    (21) appareil photographique . appareil photo
%    (65) négociant ............. marchand
%
%  DEUX PIEGES, ET CE SONT LES DEUX PLUS BEAUX DE LA COLONNE, parce
%  qu'ils ne se voient ni l'un ni l'autre : les deux mots du
%  fac-simile existent au Quebec, se comprennent au Quebec, et y
%  designent AUTRE CHOSE.
%
%    * « balayeuse » au Quebec est l'aspirateur. Le vehicule de l'ido
%      (« balayo-veturo ») est celui qui balaie la CHAUSSEE, aux cotes
%      des balayeurs (11) et dans le meme nuage de poussiere : ecrire
%      « balayeuse » aurait mis un aspirateur au milieu de la rue.
%    * « borne-fontaine » au Quebec est la bouche d'incendie. L'ido
%      ecrit « ston-fonteno », une fontaine de pierre ou l'on boit —
%      et l'alinea l'oppose justement a l'avertisseur d'incendie (76)
%      cite trois lignes plus haut. Le mot du fac-simile aurait donc
%      donne deux equipements d'incendie la ou il y en a un.
%
%  ET UN TROISIEME, DEJA RENCONTRE. « Liquoriste » retombe sur le faux
%  ami du tableau 4 : au Quebec la liqueur est une boisson gazeuse, et
%  le magasin voisin de celui du marchand de musique vend de l'alcool.
%
%  « BAS » RESTE, comme au tableau 3 : ce sont bien des bas que la
%  concierge marchande au colporteur, et c'est le mot d'ici.
% ===================================================================

%%K t14-apar-1 apar
\VUsaut{4.0mm}
\VUtitre{15.0pt}{60}{TABLEAU N\textsuperscript{o} 14}
\VUsaut{3.0mm}

%%K t14-tit-1 sub
\VUpk{11.4pt}{40}{La rue. --- Les commerçants.}
\VUsaut{3.0mm}

%%K t14-01-1 p
1. --- Mon ami Jean, qui habite la campagne, est venu passer trois jours dans ma
famille. \VUgras{Jusque-là}, il n'avait pas eu l'occasion de voir une grande
ville; aussi passons-nous toutes nos journées à nous promener, et je lui
explique ce qu'il ne connaît pas.

%%K t14-02-1 p
2. --- Nous sommes d'abord allés visiter l'\VUgras{observatoire}
\textsuperscript{(1)}, qui l'a beaucoup intéressé. En revenant par la
\VUgras{rue} \textsuperscript{(3)} où se trouve l'\VUgras{établissement de bains
turcs} \textsuperscript{(2)}, nous avons rencontré un \VUgras{enterrement}
\textsuperscript{(4)}. Devant le \VUgras{convoi funèbre} marchait le
\VUgras{clergé}, qui précédait le \VUgras{corbillard} dans lequel était placé
le \VUgras{cercueil}. Beaucoup d'amis accompagnaient la famille
\VUgras{en deuil}.

%%K t14-02-2 p
Après le passage du cortège, nous avons traversé la rue devant la grande
\VUgras{brasserie} \textsuperscript{(5)} où beaucoup de \VUgras{consommateurs}
buvaient de la \VUgras{bière} sur la terrasse, à l'abri de la \VUgras{marquise}
\textsuperscript{(6)} vitrée.

%%K t14-03-1 p
3. --- Jean était très étonné de l'\VUgras{écusson} placé entre le magasin du
\VUgras{marchand de spiritueux} \textsuperscript{(7)} et celui du
\VUgras{marchand de musique} \textsuperscript{(10)}. Il m'a demandé ce qu'il
signifiait; je lui ai répondu qu'il indiquait le \VUgras{consulat}
\textsuperscript{(8)} anglais et je lui ai fait voir le \VUgras{consul}
\textsuperscript{(9)}, qui était sur le balcon.

%%K t14-tit-2 sub
\VUpk{10.2pt}{40}{Les vitrines.}
\VUsaut{3.0mm}

%%K t14-04-1 p
4. --- Nous nous sommes bien sûr arrêtés devant l'étalage des \VUgras{instruments
de musique}, des \VUgras{harmoniums}, des \VUgras{orgues} et des
\VUgras{pianos}; nous avons regardé les couvertures illustrées des
\VUgras{partitions} et des \VUgras{morceaux} pour piano, et nous avons admiré la
\VUgras{lyre} de bois doré qui sert d'enseigne.

%%K t14-05-1 p
5. --- Les nuages de poussière soulevés par les \VUgras{balayeurs}
\textsuperscript{(11)} et la \VUgras{voiture de balayage} \textsuperscript{(12)}
nous ont obligés à passer vite devant les beaux \VUgras{meubles} du
\VUgras{marchand de meubles} \textsuperscript{(13)} et devant l'étalage de
\VUgras{poêles} et de \VUgras{lampes} du \VUgras{quincaillier}
\textsuperscript{(14)}.

%%K t14-06-1 p
6. --- À cet endroit, d'ailleurs, nous avons été obligés de quitter le trottoir,
parce qu'il était encombré de colis qu'on sortait d'une \VUgras{voiture de
déménagement} \textsuperscript{(16)} pour les porter dans un \VUgras{magasin}
\textsuperscript{(15)} qu'on venait de louer.

%%K t14-06-2 p
Cela nous a empêchés de voir les jolies voitures du \VUgras{carrossier}
\textsuperscript{(17)}, mais ne nous a pas empêchés de nous arrêter pour
regarder l'\VUgras{emballeur} \textsuperscript{(19)} qui fabriquait une caisse
pour son voisin, le \VUgras{marchand de vélos et d'automobiles}
\textsuperscript{(18)}. Ensuite, comme il était déjà un peu tard, nous sommes
rentrés à la maison.

%%K t14-tit-3 sub
\VUpk{10.2pt}{40}{Promenade en ville.}
\VUsaut{3.0mm}

%%K t14-07-1 p
7. --- Le lendemain, ma mère a proposé à Jean de le faire photographier, puis
elle m'a chargé de l'accompagner chez le \VUgras{photographe}
\textsuperscript{(20)}. Après le dîner, nous sommes allés à l'\VUgras{atelier}
de l'artiste, où nous avons dû attendre assez longtemps. Pour nous occuper, nous
avons regardé les \VUgras{portraits} \textsuperscript{(22)} accrochés au mur.

%%K t14-07-2 p
Quand notre tour est venu, le travail n'a pas été long et, dès que mon ami a eu
fini de poser devant l'\VUgras{appareil photo} \textsuperscript{(21)}, nous
sommes descendus.

%%K t14-08-1 p
8. --- Au premier étage, nous avons vu par la porte du \VUgras{coiffeur}
\textsuperscript{(23)}, qui était ouverte, son employé qui coupait les cheveux
d'un client.

%%K t14-08-2 p
Arrivés dans la rue, nous sommes passés devant le \VUgras{bureau de tabac}
\textsuperscript{(24)}. Jean a été si amusé par la \VUgras{lanterne}
\textsuperscript{(25)} rouge et par la \VUgras{grosse pipe} qui sert
d'\VUgras{enseigne} \textsuperscript{(26)} qu'il est entré en collision avec deux
vieux messieurs qui causaient \VUgras{en prenant une prise}
\textsuperscript{(27)}.

%%K t14-tit-4 sub
\VUpk{10.2pt}{40}{La pâtisserie.}
\VUsaut{3.0mm}

%%K t14-09-1 p
9. --- Les \VUgras{billets de banque} et les \VUgras{pièces de monnaie} de tous
les pays exposés dans la \VUgras{vitrine} du \VUgras{changeur}
\textsuperscript{(29)} ont attiré un moment notre attention, mais comme c'était
l'heure de la collation, nous sommes entrés chez le \VUgras{pâtissier}
\textsuperscript{(31)} sans nous arrêter plus longtemps.

%%K t14-10-1 p
10. --- Devant toutes les \VUgras{friandises} que contenait le magasin,
\VUgras{gâteaux}, \VUgras{pains d'épice}, \VUgras{biscuits} et \VUgras{bonbons}
de toutes sortes, nous avions bien du mal à choisir. Jean a fini par prendre des
\VUgras{choux à la crème}, et moi je n'ai pris qu'une petite \VUgras{tarte aux
cerises}, parce que le sucre \VUgras{me fait mal aux dents} et que je ne tiens
pas du tout à aller trop souvent chez le \VUgras{dentiste} \textsuperscript{(30)}.

%%K t14-tit-5 sub
\VUpk{10.2pt}{40}{La machine à écrire.}
\VUsaut{3.0mm}

%%K t14-11-1 p
11. --- Pour montrer à mon ami une \VUgras{machine à écrire} dont il avait
entendu parler mais qu'il ne connaissait pas, nous sommes entrés dans le
\VUgras{bureau} \textsuperscript{(32)} de mon \VUgras{cousin}
\textsuperscript{(34)}. Nous l'avons trouvé en train de dicter sa correspondance
à son \VUgras{secrétaire} \textsuperscript{(35)}, qui est \VUgras{sténographe}.
À peine une lettre était-elle finie qu'un \VUgras{dactylographe}
\textsuperscript{(33)} la recopiait à la machine. Comme nous ne voulions pas
interrompre le travail de mon parent, nous ne sommes restés chez lui que
quelques instants.

%%K t14-12-1 p
12. --- Au-dessous du bureau de mon cousin habite un grand
\VUgras{horloger-bijoutier} \textsuperscript{(37)}, qui est aussi orfèvre. Son
splendide magasin contient beaucoup d'\VUgras{horloges}, de \VUgras{pendules},
de \VUgras{montres}, de \VUgras{chaînettes} et de \VUgras{bijoux} précieux, par
exemple des \VUgras{colliers de perles}, des \VUgras{bracelets} d'or, des
\VUgras{boucles d'oreilles} et des \VUgras{bagues} ornées de \VUgras{pierres
précieuses} et de \VUgras{diamants}. Une des vitrines est spécialement réservée
à l'\VUgras{orfèvrerie} et contient une importante collection de
\VUgras{couverts} d'argent.

%%K t14-13-1 p
13. --- En tournant au coin de la rue, j'ai remarqué qu'on avait changé la
\VUgras{plaque} \textsuperscript{(36)} placée près du \VUgras{numéro} de la
maison. J'allais le dire à haute voix quand on a entendu les airs joyeux de
\VUgras{la} musique militaire. Nous nous sommes mis à courir à la rencontre du
régiment, sans penser qu'\VUgras{il} allait passer devant les magasins du
\VUgras{chapelier} \textsuperscript{(39)} et du \VUgras{gantier}
\textsuperscript{(40)}, et que nous aurions été tout aussi bien placés pour voir
son défilé en restant devant la vitrine de l'\VUgras{opticien}
\textsuperscript{(38)}, toute garnie de \VUgras{pince-nez}, de
\VUgras{lunettes} et de \VUgras{jumelles}.

%%K t14-tit-6 sub
\VUpk{10.2pt}{40}{Le défilé.}
\VUsaut{3.0mm}

%%K t14-14-1 p
14. --- En tête du \VUgras{régiment} \textsuperscript{(41)} marchait le
\VUgras{tambour-major} \textsuperscript{(42)}, suivi des \VUgras{tambours}
\textsuperscript{(43)}, des \VUgras{clairons} \textsuperscript{(44)} et de toute
la \VUgras{musique}. Le \VUgras{général} \textsuperscript{(45)}, qui était sur
la place avec son adjudant, s'était arrêté auprès du \VUgras{jet d'eau}
\textsuperscript{(49)} de la \VUgras{fontaine} \textsuperscript{(48)} pour
assister au \VUgras{défilé}.

%%K t14-14-2 p
\VUgras{Les officiers d'état-major} \textsuperscript{(46)}, le
\VUgras{colonel} et le \VUgras{lieutenant-colonel} l'ont salué de l'épée.
\VUgras{Les majors} étaient à la tête de leur \VUgras{bataillon}
\textsuperscript{(47)} et chaque \VUgras{capitaine} commandait sa
\VUgras{compagnie}, pendant que les \VUgras{lieutenants}, les
\VUgras{sous-lieutenants} et les \VUgras{sous-officiers} occupaient leur place
habituelle auprès de leurs hommes.

%%K t14-15-1 p
15. --- Beaucoup de passants s'étaient rassemblés au \VUgras{carrefour}
\textsuperscript{(50)}; les commerçants, le \VUgras{marchand de parapluies}
\textsuperscript{(51)}, le \VUgras{teinturier} \textsuperscript{(53)} et
l'\VUgras{armurier} \textsuperscript{(54)} étaient sortis pour voir les
\VUgras{soldats}. Tous les véhicules, \VUgras{fiacres} \textsuperscript{(52)},
\VUgras{voitures de maître} \textsuperscript{(57)} et même la \VUgras{tonne
d'arrosage} \textsuperscript{(56)}, s'étaient rangés le long du trottoir.

%%K t14-tit-7 sub
\VUpk{10.2pt}{40}{Scènes de la rue.}
\VUsaut{3.0mm}

%%K t14-15-2 p
Un \VUgras{représentant de commerce} \textsuperscript{(55)} qui portait à la
main ses \VUgras{boîtes d'échantillons} a traversé la rue en vitesse, et les
personnes assises sur l'\VUgras{impériale} \textsuperscript{(59)} de
l'\VUgras{omnibus} \textsuperscript{(58)} se sont retournées pour profiter du
spectacle. Un pauvre \VUgras{chanteur ambulant} \textsuperscript{(60)} avec son
\VUgras{orgue de Barbarie} \textsuperscript{(61)} a dû interrompre sa
\VUgras{chanson}, parce que le bruit des tambours couvrait sa voix.

%%K t14-16-1 p
16. --- Quand le régiment est arrivé devant le \VUgras{magasin de tissus}
\textsuperscript{(64)}, les \VUgras{couturières} \textsuperscript{(63)} et les
\VUgras{tailleurs} \textsuperscript{(62)} ont quitté leurs \VUgras{machines à
coudre} et leur travail pour regarder par la fenêtre. \VUgras{Le marchand}
\textsuperscript{(65)}, qui venait de placer de nouveaux \VUgras{costumes}
\textsuperscript{(66)} à son \VUgras{étalage}, a dû faire rentrer les pièces de
\VUgras{drap} \textsuperscript{(67)} et les \VUgras{toiles} installées sur le
trottoir, près de la \VUgras{bouche d'égout} \textsuperscript{(68)}, de peur que
la foule ne les renverse; et même un vieux \VUgras{colporteur}
\textsuperscript{(69)}, à qui une concierge marchandait des \VUgras{bas}, a été
obligé d'entrer dans un couloir pour terminer sa vente.

%%K t14-16-2 p
Nous avons accompagné le régiment jusqu'à la caserne \VUgras{et}, très contents
de notre journée, nous sommes rentrés à l'heure du souper.

%%K t14-tit-8 sub
\VUpk{10.2pt}{40}{Métiers et professions.}
\VUsaut{3.0mm}

%%K t14-17-1 p
17. --- Le lendemain, mon père nous a proposé de visiter l'imprimerie du journal
de la ville. Nous avons accepté avec enthousiasme.

%%K t14-17-2 p
L'\VUgras{imprimeur} est à la fois \VUgras{libraire} \textsuperscript{(70)}
\VUgras{(= marchand de livres)}, \VUgras{éditeur}, \VUgras{papetier} et
\VUgras{relieur}. Il a installé au premier étage la \VUgras{rédaction}
\textsuperscript{(71)} du journal, ainsi que l'\VUgras{atelier de gravure}, où
travaillent les \VUgras{lithographes} \textsuperscript{(72)} et les
\VUgras{graveurs sur cuivre} \textsuperscript{(73)}, et enfin l'\VUgras{atelier
de composition}, où se trouvent les \VUgras{typographes}
\textsuperscript{(74)}. L'\VUgras{imprimerie} \textsuperscript{(75)} proprement
dite, les \VUgras{presses} et les différentes machines sont au rez-de-chaussée.

%%K t14-tit-9 sub
\VUpk{10.2pt}{40}{Jean pose beaucoup de questions.}
\VUsaut{3.0mm}

%%K t14-18-1 p
18. --- En sortant de l'imprimerie, Jean s'est arrêté, très étonné, devant une
petite boîte vitrée placée sur une colonne, en face du magasin du
\VUgras{mercier} \textsuperscript{(77)}, et il a demandé à quoi elle servait.
Mon père lui a expliqué que c'était un \VUgras{avertisseur d'incendie}
\textsuperscript{(76)}. Tout, d'ailleurs, dans la ville, étonnait mon ami :
depuis la simple \VUgras{fontaine de pierre} \textsuperscript{(78)} et le léger
\VUgras{tricycle-porteur} \textsuperscript{(80)} conduit par un \VUgras{groom}
\textsuperscript{(79)} en livrée, jusqu'à la \VUgras{voiture des postes}
\textsuperscript{(81)}.

%%K t14-19-1 p
19. --- Pendant que mon père nous quittait un moment pour parler au
\VUgras{percepteur} \textsuperscript{(83)}, qui s'était arrêté à causer avec un
\VUgras{avocat} \textsuperscript{(82)} et un \VUgras{notaire}
\textsuperscript{(84)}, nous nous sommes amusés à suivre des yeux la circulation
dans la rue. Les personnes les plus diverses passaient devant nous : un
\VUgras{garçon pâtissier} \textsuperscript{(85)}, qui portait dans un panier un
\VUgras{pâté} splendide, venait derrière un \VUgras{marchand d'habits}
\textsuperscript{(86)}; un \VUgras{allumeur de réverbères}
\textsuperscript{(87)} causait avec un \VUgras{gazier} \textsuperscript{(88)};
une \VUgras{religieuse} \textsuperscript{(89)}, tenant une orpheline par la
main, rentrait à son couvent.

%%K t14-tit-10 sub
\VUpk{10.2pt}{40}{Sur le chemin du retour.}
\VUsaut{3.0mm}

%%K t14-20-1 p
20. --- Au moment où mon père nous a rejoints, Jean a éclaté de rire en voyant
le visage sale et l'étrange accoutrement de deux \VUgras{ramoneurs}
\textsuperscript{(91)} tout couverts de suie.

%%K t14-21-1 p
21. --- Je voulais acheter une plante à la \VUgras{fleuriste}
\textsuperscript{(92)} pour ma mère, mais mon père m'a fait remarquer qu'elle
serait bien lourde à porter et qu'il vaudrait mieux prendre des
\VUgras{oranges}. Nous nous sommes approchés de la \VUgras{marchande}
\textsuperscript{(94)} et, quand la \VUgras{gouvernante} \textsuperscript{(93)},
qui choisissait pour son élève, a eu fini ses achats, nous nous sommes fait
servir.

%%K t14-22-1 p
22. --- En rentrant à la maison, nous avons rencontré plusieurs connaissances :
une vieille amie de ma famille, appuyée au bras de sa \VUgras{dame de compagnie}
\textsuperscript{(95)}, l'\VUgras{ingénieur} \textsuperscript{(96)} des ponts et
chaussées, et une de mes cousines avec sa \VUgras{gardienne d'enfants}
\textsuperscript{(98)}.

%%K t14-22-2 p
Nous avons croisé ensuite la \VUgras{modiste} \textsuperscript{(97)} de ma mère
et deux \VUgras{moines} \textsuperscript{(99)} étrangers à la ville, qui se sont
approchés pour demander à mon père le chemin de la gare.
\VUsaut{6.0mm}
\VUfilet{22mm}
